\section{Anexo técnico}
\label{anexo:tecnico}

\subsection{Primera iteración: Redis, FAISS y LangChain}
\label{anexo:primera-iteracion-redis-faiss}

\begin{APACode}
from langchain.embeddings.openai import OpenAIEmbeddings
from langchain.vectorstores import FAISS
import os
import redis

from dotenv import load_dotenv
load_dotenv()
redis_host = os.getenv("REDIS_HOST") or "localhost"
redis_port = os.getenv("REDIS_PORT") or 6379

def connect_to_redis():
    return redis.Redis(host=redis_host, port=redis_port, db=0)

def check_index_existance(index="main") -> bool:
    redis_conn = connect_to_redis()
    index_key = f"indexes:{index}"
    if redis_conn.get(index_key) is None:
        return False
    else:
        return True

messages = ["Hello, my name is John. I am a software engineer."]
index = "main"

# FAISS
# langchain.vectorstores.faiss.FAISS
# Creamos un index general y accedemos a la data del usuario por metadata.
# {person_id:str}

embeddings = OpenAIEmbeddings()
metadatas = [{"person_id": "john", "person_name": "John"}]
conversations = FAISS.from_texts(texts=["hi"], embedding=embeddings, 
                                metadatas=metadatas)

metadatas = [{"person_id": "mary", "person_name": "Mary"}]
conversations.add_texts(texts=["hi whats up?"], metadatas=metadatas)
# ['027065ea-f01a-448d-aae2-93c05ca62768']

conversations.similarity_search(query="hi", k=2, filter={"person_id": "mary"})
# [Document(page_content='hi whats up?', 
#           metadata={'person_id': 'mary', 'person_name': 'Mary'})]
\end{APACode}
